% sections/tenancy.tex — Tenant Isolation
\section{Tenant Isolation}
\label{sec:tenancy}

\system{}'s motivating workloads come from regulated domains—audit,
compliance, safety-critical reasoning—where a single physical device
hosts the work of multiple distinct legal entities (e.g., a cloud
deployment serving several audit firms simultaneously, or one firm's
deployment serving multiple client engagements concurrently).  In
those settings, the isolation guarantee is itself a deliverable:
the operator must be able to certify, ideally by structural
argument rather than by audit log inspection, that no message
originating in one tenant's actors was ever delivered to an actor
in another tenant.  Tenancy is therefore a first-class concern of
the formal model, not a runtime checkpoint that ``usually'' holds.

This section formalizes the tenant boundary, gives a small-step
account of the cross-tenant routing rule, and proves cross-tenant
isolation: under the routing rule, no execution trace exists in
which a message crosses a tenant boundary.  The argument is by
construction: the routing function rejects cross-tenant sends as a
single first-class operation, and there is no other path by which
a message can reach an actor's inbox.  The mechanized counterpart
is in \texttt{docs/verification/tenant\_isolation.tla}.

\subsection{Audit Boundaries}
\label{sec:tenancy:boundaries}

\system{} represents the tenant of an actor as a two-tuple
\((\mathit{org\_id}, \mathit{engagement\_id}) \in \N \times \N\),
where \(\mathit{org\_id}\) identifies the legal entity (the audit
firm, the regulated organization) and \(\mathit{engagement\_id}\)
identifies the specific engagement within that entity (the audit
year, the client engagement, the compliance review).  The two-tuple
is the natural granularity for audit isolation: an audit firm
cannot legally cross-pollinate data between two of its own
engagements, even though both engagements run under the same
\(\mathit{org\_id}\).

\begin{definition}[Audit Tag]
\label{def:audit-tag}
An \emph{audit tag} is a pair \(\tau = (\mathit{org\_id},
\mathit{engagement\_id}) \in \N \times \N\).  Two tags are
\emph{equal} iff both components match.  We write \(\mathcal{T}\)
for the set of all audit tags and \(\bot \in \mathcal{T}\) for the
sentinel ``unregistered'' tag (which compares equal to nothing).
\end{definition}

Every actor carries its tag in the \(\mathit{tag}\) component of its
state vector~(\Cref{def:actor-state}); the supervisor enforces that
\(\mathit{tag}\) is set at \textsc{Spawn} and never changes thereafter.
Every K2K envelope carries the source actor's tag in its
\(\mathit{tag}\) field~(\Cref{def:k2k-envelope}).  The routing layer
checks tag equality on every \textsc{Send}.

\subsection{Routing Function}
\label{sec:tenancy:routing}

The cross-tenant rejection rule is a single inference rule that
sits between the K2K \textsc{SendOK} of \Cref{rule:k2k-send} and any
queue manipulation:

\begin{equation}
\inference[\textsc{TenantReject}]
{
  e = \langle s, d, m, \tau_e, \ldots\rangle \\
  \Theta(s) = \langle \ldots, \tau_s\rangle \\
  \Theta(d) = \langle \ldots, \tau_d\rangle \\
  \tau_s = \bot \;\lor\; \tau_d = \bot \;\lor\; \tau_s \neq \tau_d
}
{
  \langle \Theta, \langle Q, D, X\rangle, t\rangle
  \trans{\mathsf{drop}(e, \state{TenantMismatch})}
  \langle \Theta, \langle Q, D, X'\rangle, t+1\rangle
}
\label{rule:tenant-reject}
\end{equation}

\noindent where \(X' = X \cup \{(e, \state{TenantMismatch})\}\).

The corresponding accept rule has the symmetric premise:

\begin{equation}
\inference[\textsc{TenantAccept}]
{
  e = \langle s, d, m, \tau_e, \ldots\rangle \\
  \Theta(s) = \langle \ldots, \tau_s\rangle \\
  \Theta(d) = \langle \ldots, \tau_d\rangle \\
  \tau_s \neq \bot \;\land\; \tau_d \neq \bot \;\land\; \tau_s = \tau_d
}
{
  \langle \Theta, M, t\rangle
  \trans{\mathsf{tenant\_ok}(e)}
  \langle \Theta, M, t\rangle
}
\label{rule:tenant-accept}
\end{equation}

The runtime composes \textsc{TenantAccept} with \textsc{SendOK} (or
the appropriate backpressure variant): the K2K layer never sees an
envelope until the tenant check has accepted.  In the formal model,
we capture this by stating that \textsc{SendOK}~(\Cref{rule:k2k-send})
is gated by \textsc{TenantAccept}: every \textsc{SendOK} step is
preceded in the trace by a \textsc{TenantAccept} step on the same
envelope.  Failed tenant checks short-circuit to the dead-letter
store via \textsc{TenantReject}; the K2K layer has no opportunity to
see the envelope at all.

\subsection{Sub-Brokers}
\label{sec:tenancy:subbrokers}

The runtime realizes the tenant gate as a per-tag \emph{sub-broker}
inside the K2K bus.  Each sub-broker maintains its own queue
storage, its own dead-letter store, its own idempotency cache, and
its own metrics.  Sub-brokers are spawned lazily on first use of a
tag and destroyed when the tag's actors are all terminated.

The sub-broker is more than an implementation detail: it is the
mechanism by which the routing function of
\Cref{sec:tenancy:routing} becomes \emph{structural} rather than
\emph{procedural}.  An envelope in sub-broker \(\tau\) cannot reach
sub-broker \(\tau'\) because the queue storage of \(\tau'\) is in a
disjoint memory region and is reachable only through \(\tau'\)'s
public interface (\textsc{Send}, \textsc{Deliver}), which itself
checks the tag.  Even a faulty actor in \(\tau\) that attempted to
write directly into \(\tau'\)'s memory would fail at the device's
permission boundary.

In the formal model we abstract over this and treat the routing
function as a single rule pair (\textsc{TenantAccept} /
\textsc{TenantReject}); the structural argument that sub-brokers
provide an additional defense in depth is discussed in
\Cref{sec:discussion}.

\subsection{Cross-Tenant Isolation}
\label{sec:tenancy:isolation}

We can now state and prove the central isolation property.

\begin{theorem}[Cross-Tenant Isolation]
\label{thm:cross-tenant-isolation}
For every reachable configuration \(\langle \Theta, \langle Q, D,
X\rangle, t\rangle\) and every envelope \(e = \langle s, d, m,
\tau_e, \ldots\rangle \in D\), the source's audit tag and the
destination's audit tag agree:
\[
  \Theta(s).\mathit{tag} \;=\; \Theta(d).\mathit{tag}
  \quad\text{and}\quad
  \Theta(s).\mathit{tag} \neq \bot.
\]
\end{theorem}

\begin{proof}
Induction on the trace.  The base case is vacuous because \(D\) is
empty in any initial configuration.  For the induction step,
consider the rule applied at the last step: only \textsc{Deliver}
(of \Cref{rule:k2k-deliver}) adds to \(D\).  The premise of
\textsc{Deliver} requires the head of \(Q(d)\) to be \(e\), which
in turn (by the gated composition with \textsc{TenantAccept})
requires the tag check to have succeeded at the moment of
\textsc{SendOK}\((e)\).  By \textsc{TenantAccept}~(\Cref{rule:tenant-accept}),
this succeeded only if \(\Theta(s).\mathit{tag} =
\Theta(d).\mathit{tag}\) and both are non-\(\bot\).

It remains to argue that the tags have not changed between
\textsc{SendOK}\((e)\) and \textsc{Deliver}\((e)\).  By the
spawn invariant of \Cref{sec:tenancy:boundaries}, \(\mathit{tag}\)
is set at \textsc{Spawn} and never modified by any subsequent
rule.  Therefore the tag values at the time of delivery are
identical to those at the time of send, and the conclusion holds.
\qed
\end{proof}

\Cref{thm:cross-tenant-isolation} corresponds directly to
\texttt{IsolationByConstruction} in
\texttt{docs/verification/tenant\_isolation.tla} (lines 95--101).
The mechanized form keeps a stronger invariant:
\texttt{NoSameMessageInBothTenants} (lines 85--88) says that no
single message identifier appears in two different tenants' delivery
sets.  This stronger statement is implied by
\Cref{thm:cross-tenant-isolation} together with the global
uniqueness of \(\mathit{mid}\)~(see the discussion of
\Cref{def:k2k-envelope}).

\subsection{Audit Trace}
\label{sec:tenancy:audit}

The fact that \(D\) is partitioned by audit tag has an immediate
operational consequence: the audit trace of an engagement
\(\tau\) is exactly the projection of \(D\) onto envelopes whose
source has tag \(\tau\), and (by
\Cref{thm:cross-tenant-isolation}) all destinations of those
envelopes also have tag \(\tau\).  The trace is therefore
self-contained: an auditor can verify the trace of engagement
\(\tau\) without inspecting any artifact from any other engagement,
and the operator can produce that trace by querying a single
sub-broker rather than scanning the whole device.

The runtime implements this query as
\texttt{GET /api/v1/admin/k2k/dlq/list?tag=...} (for the rejected
side) and a tenant-scoped event subscription on the WebSocket
endpoint \texttt{/ws/audit/:tag} (for the delivered side).  The
operational guarantee, from the formal model, is that no envelope
emitted by another tenant will ever appear in either response.

\subsection{Limitations of the Formal Argument}
\label{sec:tenancy:limits}

The argument above proves isolation \emph{within the model}.  It
does not prove isolation against an adversary who can subvert the
runtime: an attacker who can write directly into an actor's
\(\mathit{tag}\) field would defeat the property.  We make two
notes on this:

\begin{enumerate}
\item In the runtime, \(\mathit{tag}\) is held in a memory region
that is not exposed to the actor's user code; the actor reaches its
own tag only through a read-only accessor, and the supervisor is
the sole writer.  This is a defense-in-depth measure that the
formal model abstracts over.

\item The formal model treats \textsc{Spawn} as a privileged
operation only the supervisor can invoke; the actor cannot create
new actors with arbitrary tags.  The runtime enforces this through
the supervisor RBAC layer~(\Cref{sec:implementation:admin}).
\end{enumerate}

A complete adversary model—formalizing what an attacker can do at
the GPU's hardware boundary—is beyond the scope of this paper; we
discuss it as future work in \Cref{sec:discussion}.

\subsection{Mechanization}
\label{sec:tenancy:mechanization}

The tenant-isolation rule is mechanized in
\texttt{docs/verification/tenant\_isolation.tla}, which models
\texttt{RegisterKernel} (the analogue of \textsc{Spawn} for the
isolation argument), \texttt{Send} (folded with the tenant check),
and the disjoint inboxes/rejected/delivered sets.  We model-checked
under the configuration \texttt{tenant\_isolation.cfg} for parameter
ranges \(|\textit{Tenants}| \le 2\), \(|\textit{Kernels}| \le 3\),
and \(\textit{MaxMsgs} \le 4\).  Within those bounds, the
invariants \texttt{NoSameMessageInBothTenants},
\texttt{CrossTenantRejected}, \texttt{IsolationByConstruction}, and
\texttt{InboxesWellFormed} held on every reachable state, with the
search exhausting the explored state space without violation.
